% Xiangyu Xie -- English research CV
% Portable template using standard TeX Live packages; no bundled fonts or custom class.
% Build twice: xelatex -interaction=nonstopmode -halt-on-error Xie_Xiangyu_Resume.tex
\documentclass[11pt,a4paper]{article}
\usepackage[margin=18mm,top=16mm,bottom=17mm]{geometry}
\usepackage[T1]{fontenc}
\usepackage{lmodern}
\usepackage{microtype}
\usepackage{enumitem}
\usepackage{titlesec}
\usepackage{xcolor}
\usepackage{amsmath}
\usepackage{fancyhdr}
\usepackage{hyperref}
\definecolor{ink}{HTML}{17282D}
\definecolor{accent}{HTML}{126656}
\hypersetup{colorlinks=true,urlcolor=accent,linkcolor=accent,
    pdftitle={Xiangyu Xie - Ph.D. Student at USTC - Research CV},
    pdfauthor={Xiangyu Xie},
    pdfsubject={AI for Science; Molecular Machine Learning; Scientific Agents}}
\urlstyle{same}
\color{ink}
\setlength{\parindent}{0pt}
\setlength{\parskip}{4pt}
\setlist[itemize]{leftmargin=15pt,itemsep=3pt,topsep=5pt,parsep=0pt}
\titleformat{\section}{\large\bfseries\color{accent}}{}{0pt}{}[\vspace{2pt}\titlerule]
\titlespacing*{\section}{0pt}{13pt}{7pt}
\pagestyle{fancy}
\fancyhf{}
\renewcommand{\headrulewidth}{0pt}
\fancyfoot[L]{\footnotesize Xiangyu Xie\enspace |\enspace Ph.D. Student, USTC}
\fancyfoot[R]{\footnotesize\thepage\ / 2}
\setlength{\footskip}{20pt}
\emergencystretch=2em
\newcommand{\entry}[3]{\textbf{#1}\hfill {\small #2}\\{\small #3}\par}
\begin{document}
{\fontsize{29}{32}\selectfont\bfseries Xiangyu Xie}\hfill {\small RESEARCH CV}\par
{\large Ph.D. Student at USTC}\par
Advisor: Prof. Jun Jiang\enspace |\enspace Sep 2026 -- Present\par
{\small AI for Science\enspace $\cdot$\enspace Molecular Machine Learning\enspace $\cdot$\enspace Scientific Agents}\par
{\small\href{mailto:xxy220348@mail.ustc.edu.cn}{xxy220348@mail.ustc.edu.cn}\enspace $\cdot$\enspace
\href{https://techandscixie2005.github.io/portfolio/}{Research portfolio}\enspace $\cdot$\enspace
\href{https://github.com/techandscixie2005}{GitHub: techandscixie2005}}

\section{Education}
\textbf{University of Science and Technology of China (USTC)}\par
Ph.D. Student\enspace |\enspace Sep 2026 -- Present\\
Advisor: \href{https://faculty.ustc.edu.cn/jiangjun1/en/}{Prof. Jun Jiang}\par
Undergraduate studies in Chemical Physics\enspace |\enspace From Fall 2022\\
{\small Undergraduate thesis: Fall 2025 -- Spring 2026; supervised by Prof. Jun Jiang.}

\section{Selected Research \& Projects}
\textbf{\href{https://github.com/techandscixie2005/Quantum-Agent}{QuantumAgent}}\hfill 2026\\
\textit{Evidence-Grounded Quantum Physics Teaching Agent}\\
{\small Independent Developer\enspace |\enspace Completed competition prototype}
\begin{itemize}
    \item Independently built a workflow-driven quantum physics teaching agent with a React/TypeScript workbench and a FastAPI/LangGraph backend, integrating course-grounded retrieval and scientific tools.
    \item Implemented source-traceable retrieval and a generated-code pipeline with AST checks, isolated execution, and independent numerical validation for a finite rectangular barrier case study.
    \item Developed staged learning interactions spanning initial attempts, diagnosis, teach-back, transfer, and independent problem solving, with persistent learning records and recoverable workflow state.
\end{itemize}
\vspace{5pt}
\entry{EGNN for RTP Property Prediction}{Fall 2025 -- Spring 2026}{Undergraduate Thesis\enspace |\enspace Advisor: Prof. Jun Jiang}
\begin{itemize}
    \item Built a 3D molecular graph pipeline for approximately 6,000 organic room-temperature phosphorescence (RTP) candidate molecules with quantum-chemistry labels.
    \item Implemented an eight-layer EGNN with 21-dimensional atom features and six-dimensional bond features. Compared with a character-level SMILES Transformer on the same 4,800/1,050/150 train/validation/test split.
    \item Reported test-set $R^2$ on \textbf{150 molecules}: 0.8781 for $\Delta E_{S_1-T_1}$; 0.6263 for $\mathrm{SOC}_{T_1\to S_1}$; 0.9855 for $\mathrm{SOC}_{T_1\to S_0}$; and 0.7844 for $f_{\mathrm{osc},S_1}$.
    \item Studied isolated-molecule excited-state properties; these results do not directly predict solid-state device performance, quantum yield, or phosphorescence lifetime.
\end{itemize}

\newpage
\section{Selected Research \& Projects (continued)}
\entry{Multiscale Molecular Simulation}{Fall 2024 -- Spring 2025}{Undergraduate Competition Project\enspace |\enspace Digital Intelligence Track}
\begin{itemize}
    \item Built an experiment-and-computation workflow combining GROMACS molecular dynamics with Gaussian 16 DFT calculations for n-butanol/cyclohexane systems (17:400 and 1:400).
    \item Sampled 100 trajectory frames from the dilute system for electronic-structure calculations; the project report gives a mean calculated dipole moment of 1.78 D.
    \item Used molecular environments to interpret infinite-dilution extrapolation in a physical chemistry teaching experiment; received First Prize, Central China Division, in 2025.
\end{itemize}
\vspace{4pt}
\entry{Doublet-Emissive Molecular Design}{Fall 2023 -- Fall 2024}{Undergraduate Research\enspace |\enspace Supervisor: Prof. Linsong Cui}
\begin{itemize}
    \item Designed and synthesized TTM-AnQ, coupling a TTM radical framework with a 9,10-anthraquinone acceptor through Friedel--Crafts chemistry, Suzuki coupling, and oxidation.
    \item Performed UV--Vis, fluorescence, cyclic voltammetry, and transient fluorescence characterization; observed an emission shift from 580 to 695 nm.
    \item Estimated SOMO energies of $-2.82$ and $-3.18$ eV for TTM and TTM-AnQ, respectively, supporting an exploratory study of acceptor-functionalized radical emitters.
\end{itemize}

\section{Awards and Honors}
\begin{itemize}
    \item \href{https://www.teach.ustc.edu.cn/notice/notice-info/20582.html}{\textbf{Second Prize, Graduate Division, Agent Track}} -- USTC ``107 Cup'' Computing Power and Agent Development Competition; QuantumAgent, sole team member (Sep 2026).
    \item \textbf{First Prize, Central China Division} -- 5th National Undergraduate Chemistry Experiment Innovation Design Competition, Digital Intelligence Track (2025).
    \item \textbf{USTC Outstanding Student Scholarship} (2024, 2025).
    \item \textbf{USTC Strengthening Basic Disciplines Plan Grant} (2025).
\end{itemize}

\section{Teaching Experience}
\textbf{Teaching Assistant, USTC}
\begin{itemize}
    \item \textbf{Quantum Physics} -- Spring 2026.
    \item \textbf{Mathematical Physics Equations (B)} -- Spring 2025. Facilitated problem-solving sessions and developed supplementary materials.
    \item \textbf{Complex Functions (for Chemistry)} -- Fall 2024. Led discussion sections, graded assignments, and provided one-on-one tutoring.
\end{itemize}

\section{Technical Skills}
\textbf{Molecular learning:} Python, PyTorch, EGNNs, molecular graphs, Transformer baselines.\\
\textbf{Scientific software:} React/TypeScript, FastAPI, LangGraph, PostgreSQL, Git, Linux.\\
\textbf{Chemistry \& analysis:} GROMACS, Gaussian, DFT/TD-DFT, NumPy, SciPy, spectroscopy.\\
\textbf{Research communication:} LaTeX, scientific visualization, teaching materials.
\end{document}
